\subsection{Basic Log Graphs} 
When we want to graph a logarithm, because we are graphing an \makeblank[1in], we start by choosing values for \makeblanksmall.
\begin{Example}
Sketch the graph of $f(x)=\log_2(x)$
\begin{lpic}{}{10}
	\draw (1,-2) -- (5,-2);
	\draw (3,-1) -- (3,-9);
	\regtext{2}{-2}{$x$};
	\regtext{4}{-2}{$y$};
	\foreach \x in {-3,...,3}
		\regtext{4}{-\x-6}{$\x$};
\end{lpic}
\end{Example}
What do we see in this graph?
\begin{itemize}
\item This looks like the graph of \makeblank, but \makeblank[1in] across a diagonal line.
\item The domain (as we knew) is \makeblank[1in]
\item The range (as we knew) is \makeblank[1in]
\item We can describe the end behavior, as well (though notice that one ``end'' is at \makeblank[1in]):
\begin{itemize}
\item When $x\rightarrow+\infty$, $f(x)\rightarrow\makeblanksmall$.
\item When $x\rightarrow\makeblanksmall$, $f(x)\rightarrow\makeblanksmall$.
\end{itemize}
\item We see a \makeblank[1in] asymptote, as \makeblanksmall can get close to \makeblanksmall but can never \makeblank[1in] it. In this graph, the asymptote is at \makeblank[1in].
\end{itemize}
We will see this same general shape as long as \makeblank[1in].